%!TEX root = report.tex
\documentclass[11pt,spanish]{article}


\usepackage[utf8]{inputenc}
\usepackage{subfiles}
\usepackage{biblatex}
\addbibresource{references.bib}
\usepackage{multicol}
\usepackage{amsfonts}
\usepackage{blindtext}
\usepackage{mathrsfs}
\usepackage{amsmath}
\usepackage{siunitx}
\usepackage{centernot}
\usepackage[shortlabels]{enumitem}
\usepackage{subfig}
\usepackage{datetime}
\usepackage{listingsutf8}
\usepackage[spanish]{babel}
\usepackage{tikz}
\usepackage{hyperref}
\usepackage[vlined,ruled,linesnumbered]{algorithm2e}
\usepackage{listings}
\usepackage{float}
\usepackage{url}
\usepackage{csquotes}
\usepackage{fourier} %font
\usepackage[top=2cm, bottom=2cm, left=2.5cm, right=2.5cm]{geometry}
\usepackage{pgfplots}
\usepackage{fancyhdr}
\usepackage{mdframed}
\usepackage{tikzducks}
\usepackage[nameinlink]{cleveref}
\usepackage{epigraph} 

\pgfplotsset{compat=1.18}

\usetikzlibrary{shapes.arrows, shapes.geometric, arrows.meta,angles,quotes,positioning,arrows,fit,quotes,calc}
\tikzset{>=latex} 

\setlength\algomargin{1em} 
\SetFuncSty{sc} 
\SetCommentSty{em} 


\Crefname{figure}{Fig.}{Figs.}
\newcommand\crefrangeconjunction{--}
\Crefname{table}{Tabla}{Tablas}
\Crefname{subsubsection}{Subsubsec.}{Subsubsections}
\Crefname{subsection}{Subsec.}{Subsections}
\Crefname{section}{Sec.}{Sections}
\Crefname{equation}{eq.}{eqs.}
\crefname{thm}{Theorem}{theorems}
\Crefname{thm}{Theorem}{Theorems} 

\definecolor{algoco}{rgb}{0,0.0,1}

\hypersetup{
  colorlinks=true,
  linkcolor=algoco,
  citecolor=blue,
  urlcolor=blue,
}

\lstset{
extendedchars=true
inputencoding=utf8/latin1,
basicstyle=\footnotesize\sffamily\color{black},
commentstyle=\slshape \color{gray},
numbers=left,
numbersep=10pt,
numberstyle=\tiny\color{red!80!black},
keywordstyle=\color{red!80!magenta},
showspaces=false,
showstringspaces=false,
stringstyle=\color{cyan!80!black},
tabsize=2,
literate={á}{{\'a}}1 {é}{{\'e}}1 {í}{{\'i}}1 {ó}{{\'o}}1 {ú}{{\'u}}1,
frame = single, 
numbers = none,
float, floatplacement = ht, captionpos = b,
xleftmargin = 2em, xrightmargin = 2em, 
}

\newcommand{\ub}[1]{\underbrace{#1}}
\newcommand\tcm{\textcolor{blue}}
\newcommand\tca{\textcolor{algoco}}

\setlength\epigraphwidth{.7\textwidth}
\newcommand{\tnum}{1} % reemplace 1 por el número de la tarea
\newcommand{\sem}{2026-2} % reemplace 2024-2 por el semestre correspondiente
\newcommand{\campus}{San Joaquin} % reemplace Casa Central por el campus correspondiente
\newcommand{\rolusm}{202373638-4} % reemplace 2025073100-1 por su rol
\newcommand{\namestudent}{Fernanda López Saldias} % reemplace Al Goritmo Pérez por su nombre
\newcommand{\deadline}{9 de Abril de 2026} % reemplace POR DEFINIR por la fecha de entrega

\headheight=14pt
\linespread{1.3}
\author{\namestudent}
\pagestyle{fancy}
\fancyhf{}
\fancyfoot[R]{\namestudent \\ \rolusm}
\fancyfoot[L]{Campus \campus}
\fancyfoot[C]{\thepage}
\rhead{\sem}
\lhead{INF-295}
\renewcommand{\headrulewidth}{0.4pt}
\renewcommand{\footrulewidth}{0.4pt}

\chead{ENTREGA \tnum~- PROYECTO IA}

\title{
  \huge
  \textbf{ENTREGA \tnum~: ESTADO DEL ARTE Y MODELO MATEMÁTICO} \\[1ex]
  \Large \textbf{Truck and Trailer Routing Problem (TTRP)} \\[1ex]
  \large \emph{Inteligencia Artificial --- INF-295}
}

\date{
  \small
  \today
}

\begin{document}
\maketitle
\thispagestyle{fancy}
\vspace{-1.0\baselineskip}

% --- 1. Resumen ---
\begin{abstract}
\noindent
El presente documento aborda el problema de enrutamiento de camiones y remolques (\textit{Truck and Trailer Routing Problem}, TTRP), una extensión NP-difícil del problema de enrutamiento de vehículos clásico que incorpora restricciones de accesibilidad y capacidad mediante una flota con remolques desacoplables. El objetivo principal de este informe es revisar el estado del arte de las metodologías de resolución desarrolladas en la literatura, analizando los principales enfoques exactos, heurísticos y metaheurísticos planteados desde su formulación inicial. Asimismo, se presenta la caracterización del problema, sus variantes operativas y la formulación de un modelo matemático de programación entera. Finalmente, se analizan las limitaciones de los métodos actuales y se discuten oportunidades de desarrollo para la implementación de algoritmos de optimización.
\end{abstract}
\vspace{1.5em}

\section*{Resumen}
\textit{
El presente documento aborda el problema de enrutamiento de camiones y remolques (\textit{Truck and Trailer Routing Problem}, TTRP), una extensión NP-difícil del problema de enrutamiento de vehículos clásico que incorpora restricciones de accesibilidad y capacidad mediante una flota con remolques desacoplables. El objetivo principal de este informe es revisar el estado del arte de las metodologías de resolución desarrolladas en la literatura, analizando los principales enfoques exactos, heurísticos y metaheurísticos planteados desde su formulación inicial. Asimismo, se presenta la caracterización del problema, sus variantes operativas y la formulación de un modelo matemático de programación entera. Finalmente, se analizan las limitaciones de los métodos actuales y se discuten oportunidades de desarrollo para la implementación de algoritmos de optimización.
}
\vspace{1em}

% --- 2. Introducción ---
\section{Introducción}
\input{sections/01_introduccion}

% --- 3. Definición del Problema ---
\section{Definición del Problema}
\input{sections/02_definicion_problema}

% --- 4. Estado del Arte ---
\section{Estado del Arte}
\input{sections/03_estado_del_arte}

% --- 5. Modelo Matemático ---
\section{Modelo Matemático}
\input{sections/04_modelo_matematico}

% --- 6. Conclusiones ---
\section{Conclusiones}
\input{sections/05_conclusiones}

% --- 7. Referencias Bibliográficas ---
\printbibliography[title={Referencias Bibliográficas}]

\end{document}